\subsection{The two collisions, exactly}

Both committed collisions act on the twelve lines of a cell, a line being a direction paired with its opposite, and
both are stated here in full so that the permutation claim of Section~\ref{sec:amplitude} can be read off.

\textbf{The charge rule} applies one table to the two tones $(a, b)$ meeting head-on on each line:

\begin{center}
\begin{tabular}{@{}lll@{}}
\toprule
pair & becomes & move \\
\midrule
$(-1,-1)$, $(+1,+1)$ & unchanged & like signs are inert \\
$(-1,0)$, $(0,-1)$, $(+1,0)$, $(0,+1)$ & swapped & a charge hops past a peace \\
$(0,0)$ & $(+1,-1)$ & peace creates a balanced pair \\
$(+1,-1)$ & $(-1,+1)$ & the pair flips \\
$(-1,+1)$ & $(0,0)$ & annihilation closes the cycle \\
\bottomrule
\end{tabular}
\end{center}

The table is a bijection on the nine pair states that conserves $a + b$. The create-flip-annihilate moves form a
3-cycle on the three charge-zero states, and that 3-cycle is the period-three flash of the vacuum: an all-zero
mesh becomes all $(+1,-1)$, then all $(-1,+1)$, then all zero. The rule is not an involution, so its inverse is the
inverted table, which the engine runs after the inverse stream to go backward in time.

\textbf{The momentum rule} pairs the twelve lines into six and rotates a head-on pair of equal signs into an empty
partner line: $(s, s)$ on line $i$ with $(0, 0)$ on line $j$ becomes $(0,0)$ on $i$ and $(s,s)$ on $j$, and the
reverse, for $s = \pm1$. Everything else is fixed. It is its own inverse, conserves charge, and conserves the $D_4$
momentum because the two lines of a pair are perpendicular roots. Its vacuum is a fixed point, and a lone tone never
meets the condition, which is why it streams ballistically forever.

Neither table contains a number outside $\{-1, 0, +1\}$. Composed with the stream, each is a permutation of the
$3^{24 N}$ states of an $N$-cell mesh. That is the whole of the base dynamics.
